\begin{helper}
Let \(X\) and \(W\) be measure spaces with the measures used by Lean's
\(\operatorname{volume}\), and assume that the volume measure on \(W\) is
s-finite.  Let
\(\Phi:X\to\mathbb R^3\times W\) be measure preserving, let
\(E\subseteq X\), let \(A\subseteq\mathbb R^3\times W\) be measurable, and let
\(g:\mathbb R^3\to\mathbb R\) be
continuous.  Assume \(g(x,y,z)\ge0\) whenever
\(0\le x\le y\le z\le1\).  Assume \(E\subseteq\Phi^{-1}(A)\), and assume
every point \((x,y,z,w)\in A\) satisfies \(0\le x\le y\le z\le1\).
Finally assume that, for every \(0\le x\le y\le z\le1\), the \(W\)-section
\[
\{w\in W:(x,y,z,w)\in A\}
\]
has outer measure at most \(\operatorname{ofReal}(g(x,y,z))\).  Then
\[
\operatorname{vol}(E)\le
\operatorname{ofReal}\left(\int_0^1\int_0^z\int_0^y
g(x,y,z)\,dx\,dy\,dz\right).
\]
\end{helper}
\begin{proof}
Since \(E\subseteq\Phi^{-1}(A)\), monotonicity of measure gives
\(\operatorname{vol}(E)\le \operatorname{vol}(\Phi^{-1}(A))\).  Because
\(A\) is measurable and \(\Phi\) is measure preserving,
\(\operatorname{vol}(\Phi^{-1}(A))=\operatorname{vol}(A)\).  It remains to
bound the product measure of the measurable set \(A\).

Apply the product-measure section formula successively to
\(\mathbb R\times(\mathbb R\times(\mathbb R\times W))\).  Since \(A\) is
measurable, its section \(A_{x,y,z}:=\{w:(x,y,z,w)\in A\}\) is measurable for
almost every \((x,y,z)\), and the nonnegative measurable function
\((x,y,z)\mapsto\operatorname{vol}(A_{x,y,z})\) has integral
\(\operatorname{vol}(A)\).  Equivalently,
\[
\operatorname{vol}(A)=
\int_{\mathbb R}^{+}\int_{\mathbb R}^{+}\int_{\mathbb R}^{+}
\operatorname{vol}(A_{x,y,z})\,dz\,dy\,dx
\]
with the variables ordered according to the displayed product coordinates.

The containment hypothesis says that the integrand vanishes outside the
simplex \(0\le x\le y\le z\le1\): if \(A_{x,y,z}\) is nonempty then some
\((x,y,z,w)\in A\), so the simplex inequalities hold.  Hence the preceding
integral may be restricted to this simplex.  On the simplex, the section
hypothesis gives
\(\operatorname{vol}(A_{x,y,z})\le\operatorname{ofReal}(g(x,y,z))\).
Monotonicity of the nonnegative integral therefore bounds \(\operatorname{vol}(A)\)
by the nonnegative integral of \(\operatorname{ofReal}(g(x,y,z))\) over
\(0\le x\le y\le z\le1\).

The map \((x,y,z)\mapsto g(x,y,z)\) is continuous, hence measurable and
Lebesgue integrable on the compact simplex.  The assumed nonnegativity on the
simplex allows the standard conversion between the nonnegative integral of
\(\operatorname{ofReal}(g)\) over the nested intervals
\(0\le x\le y\le z\le1\) and \(\operatorname{ofReal}\) of the corresponding
iterated real interval integral.  Thus
\[
\operatorname{vol}(A)\le
\operatorname{ofReal}\left(\int_0^1\int_0^z\int_0^y
g(x,y,z)\,dx\,dy\,dz\right).
\]
Combining this inequality with
\(\operatorname{vol}(E)\le\operatorname{vol}(A)\) proves the claim.
\end{proof}
